\documentclass[11pt]{article}
\usepackage[margin=1in]{geometry}
\usepackage{amsmath,amssymb,amsthm,booktabs}
\newtheorem{proposition}{Proposition}
\newtheorem{corollary}{Corollary}
\title{The kink without centring: Theorem 1 for a pre-norm RMSNorm block}
\author{Note for the LayerNorm-kink study (Q7a, 4 November 2026)}
\date{}
\begin{document}
\maketitle

\section{Setting}
The surrogate's first layer is $h = E z + b$ with $z = (x, u) \in \mathbb{R}^2$, $E \in \mathbb{R}^{H \times 2}$ of rank 2,
$b \in \mathbb{R}^H$, $H = 16$. LayerNorm normalises the centred vector $P h$, $P = I - \mathbf{1}\mathbf{1}^\top / H$.
RMSNorm omits the centring:
\[
  \mathrm{RMS}(h) = \frac{h}{\sqrt{\|h\|^2 / H + \varepsilon}}, \qquad
  v = \gamma \odot \mathrm{RMS}(h) + \beta .
\]
Theorem 1 was stated with $Q_E = P E$ and $c = P b$. The proposition below is the same statement with $P$ replaced by $I$.

\section{Reduction}
Let $z^\ast = -(E^\top E)^{-1} E^\top b$ (the point where $\|E z + b\|$ is smallest), $c_\perp = E z^\ast + b$ (the
part of $b$ orthogonal to the range of $E$), $D^2 = \|c_\perp\|^2 + H \varepsilon$, and $E = U \Sigma V^\top$ a thin singular
value decomposition ($U \in \mathbb{R}^{H \times 2}$).

\begin{proposition}
For every $z$, with $s = z - z^\ast$ and $w = \Sigma V^\top s = U^\top E s$,
\[
  \mathrm{RMS}(E z + b) = \sqrt{H}\, \frac{U w + c_\perp}{\sqrt{\|w\|^2 + D^2}} .
\]
\end{proposition}
\begin{proof}
$E z + b = E s + c_\perp$. Since $c_\perp \perp \mathrm{range}(E)$, $\|E z + b\|^2 = \|E s\|^2 + \|c_\perp\|^2$ and
$\|E s\| = \|w\|$. Then $\|h\|^2/H + \varepsilon = (\|w\|^2 + D^2)/H$.
\end{proof}

So the RMSNorm feature map has exactly the structure of Theorem 1: a centre $z^\ast$, a kink size $D$, half-widths
$r^\ast_x = D / \|e_x\|$ and $r^\ast_u = D / \|e_u\|$ (columns of $E$), peak feature gain $\sqrt{H}\,\sigma_{\max}(E) / D$ at
$z^\ast$, and the same far-field decay. The only change is which part of the bias sets the kink size: all of $b$ outside the
range of $E$, including its mean.

\begin{corollary}[comparison with LayerNorm for the same weights]
For the same $(E, b)$, $\|c_\perp^{\mathrm{RMS}}\| \ge \|c_\perp^{\mathrm{LN}}\|$, hence
$D_{\mathrm{RMS}} \ge D_{\mathrm{LN}}$: without centring the kink is never narrower.
\end{corollary}
\begin{proof}
$\|c_\perp^{\mathrm{LN}}\| = \min_z \|P (E z + b)\| \le \min_z \|E z + b\| = \|c_\perp^{\mathrm{RMS}}\|$ because $P$ is an
orthogonal projection.
\end{proof}

Equality needs the residual $E z^\ast_{\mathrm{RMS}} + b$ to have zero mean and the LN minimiser to coincide with it. The
difference comes from the mean of the bias (and of the encoder columns). With the encoder bias initialised at 0, both $c_\perp$
start at 0 and grow with training.

\section{Numerical check}
\texttt{theory/verify\_rms.py} (float64): over 300 random $(E, b, \varepsilon)$ with $\varepsilon \in \{0, 10^{-6}, 10^{-2}\}$ the
reduction holds to $8.9 \times 10^{-16}$ (the $U w$ form to $1.8 \times 10^{-15}$); the largest singular value of
$\partial\,\mathrm{RMS}/\partial z$ at $z^\ast$ equals $\sqrt{H}\sigma_{\max}(E)/D$ to $4 \times 10^{-9}$ (central differences);
the corollary holds on all 300 draws.

Applied to the trained LN MSE 3000 encoders (read with the RMS formula, the same $(E, b)$), the two kink sizes differ by at
most 4\,\%: $D_{\mathrm{LN}}$ from 0.132 to 0.901 and $D_{\mathrm{RMS}}$ from 0.137 to 0.908; the trained biases have mean
$|{\bar b}| \le 0.027$, so almost all of $\|b\|$ is already orthogonal to $\mathbf{1}$.

\section{What the theory does and does not predict for trained RMSNorm surrogates}
\begin{itemize}
\item It predicts the same kink structure (a centre, a size $D$, slope errors of order $1/D$) at the point where
  $\|E z + b\|$ is smallest. Since the bias starts at 0, the centre starts at the origin of $(x, u)$, as for LayerNorm.
\item For identical weights, $D_{\mathrm{RMS}} \ge D_{\mathrm{LN}}$. On LN-trained weights the difference is small, so the theory gives
  no reason to expect a qualitatively wider kink. Training an RMSNorm surrogate changes the weights, and the theory says
  nothing about how.
\item It does not predict the failure rates: these depend on the trained $D$ and on the slope errors MSE training leaves,
  as for LayerNorm.
\end{itemize}
\end{document}
